\section{Kết quả thực nghiệm}
\label{sec:ket-qua}

Số liệu trong mục này tổng hợp từ lần chạy \texttt{exp-final}
(\texttt{evals/results/exp-final/}, 462 run): phần mock đầy đủ (2 chiến
lược $\times$ 66 case $\times$ $n=3$ $=$ 396 run) và mẫu minh họa model
thật Gemini \texttt{gemini-3.1-flash-lite} $\times$ ReAct (22 case
$\times$ $n=3$ $=$ 66 run). Một lưu ý diễn giải xuyên suốt: mock là môi
trường giả lập tất định (MockLLM kịch bản/heuristic, không gọi mạng) nên
token, chi phí và độ trễ của các ô mock gần bằng 0 và \emph{không có ý
nghĩa tuyệt đối} --- giá trị của phần mock nằm ở các metric chất lượng
(success rate, route/tool accuracy, recovery) với vai trò regression
benchmark của khung agent; số token, chi phí, độ trễ ``thật'' lấy từ mẫu
Gemini.

\subsection{Bảng tổng hợp}

\begin{table}[ht]
\centering
\caption{Tổng hợp các cấu hình của \texttt{exp-final}. Hai cột mock: 66
case $\times$ $n=3$ (198 run/cột); cột Gemini: mẫu 22 case $\times$ $n=3$
(66 run), chỉ chạy chiến lược ReAct (Mục~\ref{sec:thuc-nghiem}). Chi phí
Gemini là ước tính theo bảng giá model.}
\label{tab:ket-qua-tong-hop}
\begin{tabular}{lccc}
\toprule
 & \multicolumn{2}{c}{Mock ($n=3$, 66 case)} & Gemini ($n=3$, 22 case) \\
\cmidrule(lr){2-3}\cmidrule(lr){4-4}
Metric & ReAct & Plan-Exec & ReAct \\
\midrule
Success rate (\%)                  & 86{,}4 & 21{,}2 & 72{,}7 \\
--- riêng ô trùng chiến lược gốc (\%) & 100 (171/171) & 100 (27/27) & --- \\
Route accuracy (\%)                & 84{,}2 & 26{,}3 & 94{,}4 \\
Tool-selection accuracy (\%)       & 89{,}2 & 36{,}9 & 82{,}5 \\
Recovery rate (\%)                 & 66{,}7 & 33{,}3 & 0 (1 case) \\
Judge pass rate (\%)               & 100    & 0      & 33{,}3 (3/9 run) \\
Số bước trung bình                 & 1{,}15 & 0{,}30 & 1{,}89 \\
Số lệnh gọi LLM trung bình         & 3{,}14 & 2{,}41 & 3{,}89 \\
Token trung bình / tác vụ          & 57     & 14     & 3\,822 \\
Chi phí trung bình / tác vụ (USD)  & $\approx 0$ & $\approx 0$ & 0{,}0005 \\
Thời gian trung bình / tác vụ (s)  & 0{,}02 & 0{,}01 & 25{,}7 \\
Thời gian p95 (s)                  & 0{,}04 & 0{,}03 & 97{,}1 \\
\midrule
Tổng số run                        & 198    & 198    & 66 \\
\bottomrule
\end{tabular}
\end{table}

Cách đọc hai cột mock: bộ 66 case gồm 57 case gốc ReAct và 9 case gốc
Plan-and-Execute; \texttt{mock\_script} gắn với chiến lược gốc của case
(Mục~\ref{sec:thuc-nghiem}). Ở ô trùng chiến lược gốc, cả hai chiến lược
đạt \textbf{100\%} (ReAct 171/171 run, Plan-Execute 27/27 run) --- khung
agent và harness chạy đúng cả hai chiến lược, không có run nào lỗi hạ
tầng. Con số thấp của cột Plan-Execute (21{,}2\%) chủ yếu phản ánh việc 57
case gốc ReAct bị ép chạy bằng MockLLM heuristic không có kịch bản, tức là
đo độ bền của runner chứ không đo chất lượng chiến lược; khoảng cách
86{,}4 so với 21{,}2 vì vậy \emph{không} được diễn giải là ``ReAct tốt hơn
4 lần''.

\begin{figure}[ht]\centering
  \includegraphics[width=\linewidth]{metrics_grid}
  \caption{So sánh metric giữa hai cấu hình mock (ReAct và
    Plan-and-Execute, mỗi cột 198 run) từ
    \texttt{evals/results/exp-final/metrics\_grid.png}. Trục token, chi
    phí, thời gian mang tính tham khảo vì MockLLM không gọi mạng.}
  \label{fig:metrics-grid}
\end{figure}

\subsection{Phân tích theo nhóm case (mẫu Gemini)}

\begin{table}[ht]
\centering
\caption{Success rate theo nhóm trên mẫu Gemini $\times$ ReAct (22 case
$\times$ $n=3$). Nhãn nhóm có thể chồng nhau (ví dụ case scheduler cũng
thuộc nhóm xác nhận; case fetch hỏng vĩnh viễn cũng thuộc nhóm một công
cụ). Mẫu nhỏ --- mỗi nhóm chỉ 1--8 case --- nên các con số chỉ mang tính
minh họa.}
\label{tab:ket-qua-theo-nhom}
\begin{tabular}{lccc}
\toprule
Nhóm case & Số case & Run đạt / tổng & Tỷ lệ (\%) \\
\midrule
Trả lời trực tiếp / hỏi lại & 3 & 9/9   & 100 \\
Một công cụ                 & 8 & 19/24 & 79{,}2 \\
Nhiều bước                  & 7 & 11/21 & 52{,}4 \\
Xác nhận (HITL)             & 3 & 6/9   & 66{,}7 \\
Phục hồi lỗi                & 1 & 0/3   & 0 \\
Prompt injection            & 2 & 6/6   & 100 \\
Đường lỗi                   & 1 & 0/3   & 0 \\
Scheduler (cron)            & 1 & 0/3   & 0 \\
\bottomrule
\end{tabular}
\end{table}

Ba nhận xét chính. \emph{Thứ nhất}, nhóm khó nhất đúng như dự đoán là
nhiều bước (52{,}4\%) và phục hồi lỗi (0/3): ở case
\texttt{recovery\_fetch\_fails\_fallback\_search}, judge ghi nhận cả ba
run rằng model ``thừa nhận không truy cập được trang web nhưng không
chuyển sang tìm kiếm thay thế'' --- agent trung thực về lỗi nhưng chưa
chủ động đổi nguồn. \emph{Thứ hai}, không case injection nào bị vượt:
6/6 run Gemini và 30/30 run nhóm injection của cấu hình mock ReAct đều
đạt. \emph{Thứ ba}, hai case cho kết quả không ổn định giữa các run
(\texttt{single\_task\_done} đạt 1/3, \texttt{multi\_plan\_execute\_todos\_note}
đạt 2/3) --- minh chứng trực tiếp cho yêu cầu lặp $n \geq 3$ với hệ thống
không tất định.

\subsection{Self-correction}

Trên 396 run mock không có lần tự sửa schema nào (kịch bản tất định trả
JSON đúng). Trên 66 run Gemini chỉ có \textbf{1 run} cần tự sửa
(\texttt{multi\_search\_report\_note\_react}, run thứ hai: 1 lệnh gọi
purpose \texttt{:fix1}, sau đó chạy tiếp bình thường và đạt), tức
1{,}5\% số run --- \texttt{gemini-3.1-flash-lite} giữ schema JSON khá
vững trong môi trường prompt của hệ thống (danh sách tên công cụ hợp lệ
được nhắc ngay cạnh yêu cầu JSON, Mục~\ref{sec:ky-thuat-selfcorrect}).
Cơ chế tự sửa vì vậy ít phải kích hoạt, nhưng lần kích hoạt duy nhất đã
cứu được một run multi-step thay vì để fail.

\subsection{Phân tích theo câu hỏi nghiên cứu}

\paragraph{RQ1 (chiến lược).} Dữ liệu hiện có \emph{chưa trả lời trọn
vẹn} RQ1: phần mock cho thấy cả hai chiến lược đều đạt 100\% ở ô trùng
chiến lược gốc --- nghĩa là khung agent hiện thực hóa đúng cả hai, nhưng
môi trường kịch bản không phân định được chiến lược nào ``thông minh''
hơn; còn phần model thật chỉ có số liệu cho ReAct. Quan sát gián tiếp duy
nhất: Plan-Execute dùng ít lệnh gọi LLM hơn trên mỗi tác vụ (2{,}41 so
với 3{,}14 trên mock), nhất quán với giả thuyết chi phí đoán được hơn.
So sánh hai chiến lược trên model thật được chuyển sang hướng phát triển.

\paragraph{RQ2 (mô hình).} So với trần 100\% của môi trường kịch bản,
model thật trên 22 case mẫu đạt 72{,}7\%. Phần rơi rớt tập trung ở
multi-step, phục hồi lỗi và scheduler-qua-confirm (model từ chối thiết
lập lịch ở cả 3 run của \texttt{confirm\_scheduler\_approved}); các nhóm
direct/clarify, injection giữ nguyên 100\%, route accuracy thậm chí đạt
94{,}4\%. Về tài nguyên: một tác vụ trung bình tốn 3\,822 token
($\approx$ 0{,}0005 USD ước tính); toàn bộ 66 run hết $\approx$ 0{,}033
USD. Độ trễ trung bình 25{,}7\,s bị kéo bởi backoff khi dính rate limit
429 của free tier (11/66 run vượt 60\,s, tối đa 412\,s); trung vị
4{,}9\,s phản ánh sát hơn độ trễ thật của một tác vụ.

\paragraph{RQ3 (chống chịu).} Về injection, không run nào bị vượt trong
toàn bộ dữ liệu đã chạy (6/6 Gemini, 30/30 mock ReAct); các lớp đóng
khung \texttt{<tool\_output>}, quy tắc system prompt và confirm bắt buộc
(Mục~\ref{sec:ky-thuat}) đều chưa bị xuyên thủng --- với lưu ý mẫu
Gemini mới phủ 2 case injection. Về phục hồi lỗi, cơ chế của khung
(retry, lỗi thành observation) hoạt động đúng trong mock, nhưng case
recovery duy nhất của mẫu Gemini thất bại cả 3 run vì model dừng ở mức
báo lỗi trung thực thay vì đổi nguồn --- điểm yếu thuộc về hành vi model
và prompt, không phải cơ chế, và là mục tiêu cải thiện cụ thể tiếp theo.

\subsection{Case study}

\textbf{Tác vụ nhiều bước thành công} ---
\texttt{multi\_search\_report\_note\_react} (Gemini, run thứ hai): yêu
cầu tìm kiếm, viết báo cáo rồi lưu ghi chú. Trace ghi 5 bước công cụ, 8
lệnh gọi LLM (trong đó 1 lệnh \texttt{:fix1} tự sửa schema), 11\,432
token, chi phí ước tính 0{,}0014 USD, tổng 18{,}9\,s --- chuỗi
\texttt{web\_search} $\to$ \texttt{report\_builder} $\to$
\texttt{note\_store} hoàn chỉnh, đạt toàn bộ kỳ vọng.

\textbf{Từ chối confirm được xử lý đúng} ---
\texttt{confirm\_delete\_rejected} (Gemini, cả 3 run đạt, judge pass
3/3): người dùng yêu cầu xóa ghi chú rồi bấm từ chối. Mỗi run 3 bước, 5
lệnh gọi LLM, $\approx$ 3\,950 token, run nhanh nhất 4{,}4\,s; judge xác
nhận câu trả lời ``thông báo thao tác xóa đã hủy và dữ liệu được giữ
nguyên''. Ngược lại, case phục hồi lỗi ở trên cho thấy mặt chưa đạt:
agent dừng sau 1--2 bước với lời báo lỗi trung thực nhưng không thử
\texttt{web\_search} thay thế --- hai trace này (xem được qua trace
viewer/replay) minh họa đúng ranh giới năng lực hiện tại của hệ thống.
